\documentclass[12pt]{article}
\usepackage[utf8]{inputenc}
\usepackage[T1]{fontenc}
\usepackage{hyperref}
\title{Western Subanon Language Resources for K-12 Curriculum Development}
\author{Compiled by Claude Sonnet 5.5 (claude-sonnet-5-5) for Andr\'as Kornai}
\date{}
\begin{document}
\maketitle

\section{Introduction}
Western Subanon (Sinubanon; ISO 639-3: suc; Glottocode: west2557) is an Austronesian language of the Zamboanga Peninsula, Philippines. Speaker estimates range from about 120,000 to 300,000. This report rests on a single web search.

\section{Primary Educational Resources}
No ready-made school textbooks or mother-tongue curriculum were found in the search results. Omniglot states that a dictionary and literacy materials are being produced, but I could not confirm that they have been published.
\begin{itemize}
\item Omniglot, Western Subanon (Sinubanon): \url{https://omniglot.com/writing/westernsubanon.htm}
\end{itemize}

\section{Linguistic and Reference Materials}
\begin{itemize}
\item \emph{A Grammar of Western Subanon} (dissertation; De La Salle University repository record; the author attribution differs between records): \url{https://animorepository.dlsu.edu.ph/faculty_research/9373}
\item Featured talk at the 2020 LLL Conference on ``A Grammar of Western Subanon'' (University of Hawai\textquoteleft i ScholarSpace): \url{https://scholarspace.manoa.hawaii.edu/items/fdca16c7-9225-4a10-861a-5fb4dbb43d05}
\item \emph{Western Subanon Grammar} (Brill, 2025, open access; Sharon Bulalang): \url{https://www.degruyterbrill.com/document/isbn/9789004703056/html}
\item Wikipedia, Subanen languages: \url{https://en.wikipedia.org/wiki/Subanen_languages}
\end{itemize}

\section{Religious and Community Resources}
\begin{itemize}
\item Joshua Project, Western Kalibugan Subanon: \url{https://joshuaproject.net/people_groups/15100}
\end{itemize}

\section{Digital and Audio Resources}
The documentary corpus accompanying the grammar is reported to be in the Kaipuleohone Digital Ethnographic Archive at the University of Hawai\textquoteleft i. No direct archive URL appeared in the search results.

\section{Curriculum Development Context}
The language is used as a language of instruction in some schools. The recent open-access grammar and its corpus are a strong basis for materials. The grammar's author, a Western Subanon speaker and former grade-school teacher, states a goal of supporting curriculum development.
\end{document}